% from ega3-1-fr.tex
\oldpage[0\textsubscript{III}]{41}
\[
\label{0.11.6.4.1-fr}
\tag{11.6.4.1}
{}'E^2_{pq}=H_p(L_qT(K_\bullet,K'_\bullet))
\]
\[
\label{0.11.6.4.2-fr}
\tag{11.6.4.2}
{}''E^2_{pq}=
\sum_{q'+q''=q}L_p(H_{q'}(K_\bullet),H_{q''}(K'_\bullet)).
\]
Ici encore, c'est la \emph{seconde} suite spectrale qui est régulière, les deux
suites étant birégulières lorsqu'il s'agit de complexes $K_\bullet$,
$K'_\bullet$ limités inférieurement, ou quand les objets des catégories
abéliennes que l'on considère ont des résolutions projectives de longueur
fixée.

En outre :
\end{env}

\begin{proposition}[11.6.5]
\label{0.11.6.5-fr}
Soient $\boldsymbol{C}$, $\boldsymbol{C}'$, $\boldsymbol{C}''$ trois catégories
abéliennes, $T$ un bifoncteur covariant biadditif de
$\boldsymbol{C}\times\boldsymbol{C}'$ dans $\boldsymbol{C}''$.
\begin{enumerate}
\item[{\rm(i)}] {$\mathbf{L}_\bullet T(K_\bullet,K'_\bullet)$ est un bifoncteur
homologique en chacun des complexes limités inférieurement $K_\bullet$,
$K'_\bullet$ (formés respectivement d'objets de $\boldsymbol{C}$ et d'objets de
$\boldsymbol{C}'$).}
\item[{\rm(ii)}] {Supposons $K_\bullet$ et $K'_\bullet$ limités
inférieurement. Soient $L_{\bullet\bullet}=(L_{ij})$,
$L'_{\bullet\bullet}=(L'_{ij})$ deux bicomplexes tels que $L_{ij}=0$ et
$L'_{ij}=0$ pour $j<0$, que pour tout $i$, $L_{i,\bullet}$ soit une
résolutiou de $K_i$ et $L'_{i,\bullet}$ une résolution de $K'_i$, et enfin que
$L_nT(L_{ij},L'_{hk})=0$ pour $n>0$ et tout système $(i,j,h,k)$. Alors on a
un isomorphisme fonctoriel}
\[
\label{0.11.6.5.1-fr}
\tag{11.6.5.1}
\mathbf{L}_\bullet T(K_\bullet,K'_\bullet)
\xrightarrow{\sim} H_\bullet(T(L_{\bullet\bullet},L'_{\bullet\bullet})).
\]
\item[{\rm(iii)}] {Supposons en outre que pour tout couple $(i,j)$ et
tout couple $(h,k)$, les foncteurs
$A\leadsto T(A,L'_{hk})$ et $A'\leadsto T(L_{ij},A')$
soient exacts dans $\boldsymbol{C}$ et $\boldsymbol{C}'$ respectivement. Alors on a
des isomorphismes fonctoriels}
\[
\label{0.11.6.5.2-fr}
\tag{11.6.5.2}
\mathbf{L}_\bullet T(K_\bullet,K'_\bullet)
\xrightarrow{\sim} H_\bullet(T(L_{\bullet\bullet},K'_\bullet))
\xrightarrow{\sim} H_\bullet(T(K_\bullet,L'_{\bullet\bullet})).
\]
\end{enumerate}
\upshape Les démonstrations sont analogues à celles de
(\hyperref[0.11.5.5-fr]{11.5.5}) et (\hyperref[0.11.5.6-fr]{11.5.6}).
\end{proposition}

\subsection{Hyperhomologie d'un foncteur par rapport à un bicomplexe
$K_{\bullet\bullet}$}
\label{subsection:0.11.7-fr}

\begin{env}[11.7.1]
\label{0.11.7.1-fr}
Soit $\boldsymbol{C}$ une catégorie abélienne dans laquelle tout objet est
quotient d'un objet projectif. Considérons un bicomplexe
$K_{\bullet\bullet}=(K_{ij})$ formé d'objets de $\boldsymbol{C}$, et dont les deux
degrés sont \emph{limités inférieurement}; on peut toujours se borner au cas où
$K_{ij}=0$ pour $i<0$ ou $j<0$, et c'est ce que nous ferons désormais. On
peut considérer $K_{\bullet\bullet}$ comme un complexe (simple) formé
d'objets $K_{i,\bullet}=(K_{ij})_{j\geq 0}$ de la catégorie abélienne
$\boldsymbol{K}$ des complexes à degrés positifs d'objets de $\boldsymbol{C}$. Il
résulte du lemme (\hyperref[0.11.5.2.1-fr]{11.5.2.1}) (ou plutôt du lemme
«~dual~» obtenu par «~renversement des flèches~») et de (M, V, 2.2) que
$K_{\bullet\bullet}$ admet une \emph{résolution projective de
Cartan-Eilenberg} dans la catégorie $\boldsymbol{K}$; une telle résolution est un
\emph{tricomplexe} $M_{\bullet\bullet\bullet}=(M_{ijk})$ de $\boldsymbol{C}$, à degrés
tous $\geq 0$, formé d'objets projectifs, tel que pour tout $i$,
$M_{i,\bullet\bullet}$, $\mathrm B_i^{\mathrm I}(M_{\bullet\bullet\bullet})$,
$\mathrm Z_i^{\mathrm I}(M_{\bullet\bullet\bullet})$,
$H_i^{\mathrm I}(M_{\bullet\bullet\bullet})$ constituent des \emph{résolutions
projectives} de $K_{i,\bullet}$,
$\mathrm B_i^{\mathrm I}(K_{\bullet\bullet})$,
$\mathrm Z_i^{\mathrm I}(K_{\bullet\bullet})$,
$H_i^{\mathrm I}(K_{\bullet\bullet})$ respectivement dans la catégorie
$\boldsymbol{K}$; en particulier, pour tout couple $(i,j)$,
$M_{ij,\bullet}$ est une \emph{résolution projective} de $K_{ij}$ dans
$\boldsymbol{C}$.
\end{env}

\begin{proposition}[11.7.2]
\label{0.11.7.2-fr}
Soit $T$ un foncteur covariant additif de $\boldsymbol{C}$ dans une catégorie
abélienne $\boldsymbol{C}'$. Avec les notations de
\textup{(\hyperref[0.11.7.1-fr]{11.7.1})}, {l'homologie
$H_\bullet(T(M_{\bullet\bullet\bullet}))$ du complexe simple associé au
tricomplexe $T(M_{\bullet\bullet\bullet})$ est canoniquement isomorphe à
l'hyperhomologie $\mathbf{L}_\bullet T(K_{\bullet\bullet})$ du complexe
